\chapter{2004年北京卷（秋文）}
\section{单选题}
\begin{problem}
设 \(M=\{x\mid -2\leqslant x\leqslant 2\}\)，\(N=\{x\mid x<1\}\)，则 \(M\cap N\) 等于（\quad）

\choices
    {\(\{x\mid 1<x<2\}\)}
    {\(\{x\mid -2<x<1\}\)}
    {\(\{x\mid 1<x\leqslant 2\}\)}
    {\(\{x\mid -2\leqslant x<1\}\)}
\begin{answer}
D
\end{answer}

\begin{solution}
因为 \(M\) 表示区间 \([-2,2]\)，而 \(N\) 表示区间 \((-\infty,1)\)，所以
\(M\cap N=\{x\mid -2\leqslant x<1\}\)，正确选项为 D.
\end{solution}

\end{problem}

\begin{problem}
满足条件 \(\left|z\right|=\left|3+4\i\right|\) 的复数 \(z\) 在复平面上对应点的轨迹是（\quad）

\choices
    {一条直线}
    {两条直线}
    {圆}
    {椭圆}
\begin{answer}
C
\end{answer}

\begin{solution}
设 \(z=x+y\i\)，则
\(\left\lvert 3+4\i\right\rvert=\sqrt{3^{2}+4^{2}}=5\)，从而
\(\left\lvert z\right\rvert=5\) 等价于 \(x^{2}+y^{2}=25\). 这表示以原点为圆心、半径为 \(5\) 的圆，故正确选项为 C.
\end{solution}

\end{problem}

\begin{problem}
设 \(m\)，\(n\) 是两条不同的直线，\(\alpha\)，\(\beta\)，\(\gamma\) 是三个不同的平面，给出下列四个命题：

\begin{circlelist}
\item 若 \(m\perp\alpha\)，\(n\parallel\alpha\)，则 \(m\perp n\)；
\item 若 \(\alpha\parallel\beta\)，\(\beta\parallel\gamma\)，\(m\perp\alpha\)，则 \(m\perp\gamma\)；
\item 若 \(m\parallel\alpha\)，\(n\parallel\alpha\)，则 \(m\parallel n\)；
\item 若 \(\alpha\perp\gamma\)，\(\beta\perp\gamma\)，则 \(\alpha\parallel\beta\).
\end{circlelist}

其中正确命题的序号是（\quad）

\choices
    {\circled{1}和\circled{2}}
    {\circled{2}和\circled{3}}
    {\circled{3}和\circled{4}}
    {\circled{1}和\circled{4}}
\begin{answer}
A
\end{answer}

\begin{solution}
若 \(m\perp\alpha\)，则 \(m\) 垂直于平面 \(\alpha\) 内的任意直线；而 \(n\parallel\alpha\) 表示 \(n\) 的方向与平面 \(\alpha\) 内某条直线的方向相同，因此 \(m\perp n\)，第一个命题正确.

由 \(\alpha\parallel\beta\)，\(\beta\parallel\gamma\) 及三个平面互不重合，得 \(\alpha\parallel\gamma\). 平行平面的法向方向相同，所以 \(m\perp\alpha\) 时也有 \(m\perp\gamma\)，第二个命题正确.

第三个命题不正确：在一个与 \(\alpha\) 平行的平面内取两条相交直线 \(m\)，\(n\)，它们都平行于 \(\alpha\)，但彼此不平行. 第四个命题也不正确：两个互相垂直于同一平面的平面可以相交，例如两个不同的竖直平面都垂直于水平面. 因而正确命题为第一和第二个，选 A.
\end{solution}

\end{problem}

\begin{problem}
已知 \(a\)，\(b\)，\(c\) 满足 \(c<b<a\) 且 \(ac<0\)，那么下列选项中一定成立的是（\quad）

\choices
    {\(ab>ac\)}
    {\(c(b-a)<0\)}
    {\(cb^{2}<ab^{2}\)}
    {\(ac(a-c)>0\)}
\begin{answer}
A
\end{answer}

\begin{solution}
由 \(c<b<a\) 且 \(ac<0\) 可知 \(a\) 与 \(c\) 异号. 若 \(a<0\)，则 \(c<b<a<0\)，从而 \(c<0\)，这将导致 \(ac>0\)，矛盾. 所以 \(a>0\)，进而 \(c<0\).

因为 \(b>c\) 且 \(a>0\)，所以 \(ab>ac\)，选项 A 一定成立. 另外，\(c<0\)，\(b-a<0\)，故 \(c(b-a)>0\)，B 不成立；由 \(b^{2}\geqslant0\) 只能得到 \(cb^{2}\leqslant ab^{2}\)，当 \(b=0\) 时等号成立，C 不一定成立；最后 \(ac<0\)，\(a-c>0\)，所以 \(ac(a-c)<0\)，D 不成立. 因此选 A.
\end{solution}

\end{problem}

\begin{problem}
从长度分别为 \(1\)，\(2\)，\(3\)，\(4\)，\(5\) 的五条线段中，任取三条的不同取法共有 \(n\) 种. 在这些取法中，以取出的三条线段为边可组成的钝角三角形的个数为 \(m\)，则 \(\frac{m}{n}\) 等于（\quad）

\choices
    {\(\frac{1}{10}\)}
    {\(\frac{1}{5}\)}
    {\(\frac{3}{10}\)}
    {\(\frac{2}{5}\)}
\begin{answer}
B
\end{answer}

\begin{solution}
任取三条线段的取法数为
\(n=\mathrm C_{5}^{3}=10\). 将所有取法按边长从小到大列出：\((1,2,3)\)、\((1,2,4)\)、\((1,2,5)\)、\((1,3,4)\)、\((1,3,5)\)、\((1,4,5)\)、\((2,3,4)\)、\((2,3,5)\)、\((2,4,5)\)、\((3,4,5)\).

三角形的判定条件是两条较短边之和大于最长边. 因而 \((1,2,3)\)、\((1,3,4)\)、\((1,4,5)\)、\((2,3,5)\) 的两短边之和等于最长边，不能组成三角形；\((1,2,4)\)、\((1,2,5)\)、\((1,3,5)\) 的两短边之和小于最长边，也不能组成三角形. 只有 \((2,3,4)\)、\((2,4,5)\)、\((3,4,5)\) 能组成三角形.

对于边长依次为 \(a\leqslant b\leqslant c\) 的三角形，由余弦定理，最大角为钝角、直角或锐角，分别对应 \(c^{2}>a^{2}+b^{2}\)、\(c^{2}=a^{2}+b^{2}\) 或 \(c^{2}<a^{2}+b^{2}\). 对 \((2,3,4)\)，有 \(4^{2}=16>13=2^{2}+3^{2}\)；对 \((2,4,5)\)，有 \(5^{2}=25>20=2^{2}+4^{2}\)，这两组是钝角三角形. 对 \((3,4,5)\)，有 \(5^{2}=3^{2}+4^{2}\)，它是直角三角形. 因而 \(m=2\)，所以
\(\dfrac{m}{n}=\dfrac{2}{10}=\dfrac15\)，选 B.
\end{solution}

\end{problem}

\begin{problem}
如图，在正方体 \(ABCD-A_{1}B_{1}C_{1}D_{1}\) 中，\(P\) 是侧面 \(BB_{1}C_{1}C\) 内一动点，若 \(P\) 到直线 \(BC\) 与直线 \(C_{1}D_{1}\) 的距离相等，则动点 \(P\) 的轨迹所在的曲线是（\quad）

\bitmapfigure[width=0.40\linewidth]{img/2004/beijing_liberal_autumn/q6_fig1.png}

\choices
    {直线}
    {圆}
    {双曲线}
    {抛物线}
\begin{answer}
D
\end{answer}

\begin{solution}
设正方体棱长为 \(l>0\)，建立直角坐标系，使
\(A=(0,0,0)\)，\(B=(l,0,0)\)，\(C=(l,l,0)\)，\(C_{1}=(l,l,l)\)，且侧面 \(BB_{1}C_{1}C\) 为平面 \(x=l\). 设 \(P=(l,y,z)\)，其中 \(0\leqslant y\)，\(z\leqslant l\).

点 \(P\) 到直线 \(BC\) 的距离为 \(z\). 直线 \(C_{1}D_{1}\) 的点可写为 \((s,l,l)\)，所以
\[
\operatorname{dist}(P,C_{1}D_{1})^{2}=(l-y)^{2}+(l-z)^{2}
\]
其中最近点对应于 \(s=l\)，即 \(C_{1}\). 因此题设条件等价于
\[
z=\sqrt{(l-y)^{2}+(l-z)^{2}}
\]
这正是平面 \(x=l\) 内的点到焦点 \(C_{1}\) 的距离等于到准线 \(BC\) 的距离的轨迹，即抛物线在该正方形侧面内的部分. 故选 D.
\end{solution}

\end{problem}

\begin{problem}
函数 \(f(x)=x^{2}-2ax-3\) 在区间 \([1,2]\) 上存在反函数的充分必要条件是（\quad）

\choices
    {\(a\in(-\infty,1]\)}
    {\(a\in[2,+\infty)\)}
    {\(a\in(-\infty,1]\cup[2,+\infty)\)}
    {\(a\in[1,2]\)}
\begin{answer}
C
\end{answer}

\begin{solution}
函数可写为
\(f(x)=(x-a)^{2}-a^{2}-3\)，其对称轴为 \(x=a\). 二次函数在闭区间 \([1,2]\) 上存在反函数，当且仅当它在该区间上具有单调性，也就是对称轴不在区间内部.

当 \(a\leqslant1\) 时，\(f'(x)=2(x-a)\geqslant0\)（除端点外为正），函数在 \([1,2]\) 上严格递增；当 \(a\geqslant2\) 时，\(f'(x)\leqslant0\)（除端点外为负），函数在 \([1,2]\) 上严格递减. 若 \(1<a<2\)，取足够小的 \(t>0\)，使 \(a-t\)，\(a+t\in[1,2]\)，则
\(f(a-t)=f(a+t)\)，函数不具有单调性，也不可能有反函数. 所以
\(a\in(-\infty,1]\cup[2,+\infty)\)，选 C.
\end{solution}

\end{problem}

\begin{problem}
函数 \(f(x)=\begin{cases}
x,&x\in P,\\
-x,&x\in M,
\end{cases}\) 其中 \(P\)，\(M\) 为实数集 \(\R\) 的两个非空子集，又规定 \(f(P)=\{y\mid y=f(x),x\in P\}\)，\(f(M)=\{y\mid y=f(x),x\in M\}\)，给出下列四个判断：

\begin{circlelist}
\item 若 \(P\cap M=\varnothing\)，则 \(f(P)\cup f(M)=\varnothing\)；
\item 若 \(P\cap M\neq\varnothing\)，则 \(f(P)\cap f(M)\neq\varnothing\)；
\item 若 \(P\cup M=\R\)，则 \(f(P)\cup f(M)=\R\)；
\item 若 \(P\cup M\neq\R\)，则 \(f(P)\cup f(M)\neq\R\).
\end{circlelist}

其中正确判断有（\quad）

\choices
    {3个}
    {2个}
    {1个}
    {0个}
\begin{answer}
B
\end{answer}

\begin{solution}
因为题目给出的 \(f\) 是函数，若 \(x\in P\cap M\)，则同一个 \(x\) 必须满足 \(f(x)=x\) 和 \(f(x)=-x\)，故 \(x=0\)，即 \(P\cap M\subseteq\{0\}\). 对任意 \(p\in P\)，即使 \(p\in M\)，也只能有 \(p=0\)，仍有 \(f(p)=p\)，所以 \(f(P)=P\)；同理 \(f(M)=-M=\{-m\mid m\in M\}\).

第一个命题不正确：\(P\)，\(M\) 都是非空集，即使 \(P\cap M=\varnothing\)，\(f(P)\) 和 \(f(M)\) 仍分别非空，所以它们的并集不可能为空.

第二个命题正确：若 \(P\cap M\neq\varnothing\)，由上面的结论必有 \(0\in P\cap M\)，而 \(f(0)=0\)，所以 \(0\in f(P)\cap f(M)\).

第三个命题不正确. 取 \(P=[0,+\infty)\)，\(M=(-\infty,0]\)，则 \(P\cup M=\R\)，且
\(f(P)=[0,+\infty)\)，\(f(M)=[0,+\infty)\)，故 \(f(P)\cup f(M)\neq\R\).

证明第四个命题的逆否命题. 假设 \(f(P)\cup f(M)=\R\)，若存在 \(x\notin P\cup M\)，则 \(x\notin P\)，从而 \(x\notin f(P)=P\)，所以 \(x\in f(M)\). 由 \(f(M)=-M\) 可知 \(-x\in M\).

又因 \(-x\in\R=f(P)\cup f(M)=P\cup(-M)\)，若 \(-x\in f(M)=-M\)，则存在 \(m\in M\) 使 \(-m=-x\)，即 \(m=x\)，与 \(x\notin M\) 矛盾. 因而 \(-x\in f(P)=P\). 结合 \(-x\in M\)，得 \(-x\in P\cap M\subseteq\{0\}\)，所以 \(x=0\). 但此时 \(x=0\in M\)，仍与 \(x\notin P\cup M\) 矛盾. 所以 \(P\cup M=\R\)，第四个命题正确.

因此正确判断有第二和第四个，共 \(2\) 个，选 B.
\end{solution}

\end{problem}

\section{填空题}
\begin{problem}
函数 \(y=\sin x\cos x\) 的最小正周期是\fillinblank{}.
\begin{answer}
\(\pi\)
\end{answer}

\begin{solution}
利用二倍角公式，\(y=\sin x\cos x=\dfrac12\sin2x\). 因为 \(\sin 2(x+\pi)=\sin(2x+2\pi)=\sin2x\)，所以 \(\pi\) 是该函数的正周期. 若正数 \(T\) 是它的周期，则对任意 \(x\)，有 \(\sin(2x+2T)=\sin2x\). 取 \(x=0\) 和 \(x=\dfrac{\pi}{4}\)，分别得到 \(\sin2T=0\) 以及 \(\cos2T=1\)，故 \(2T=2k\pi\)（\(k\in\Z\)），从而 \(T=k\pi\). 正数周期中最小者为 \(\pi\).
\end{solution}

\end{problem}

\begin{problem}
方程 \(\lg(x^{2}+2)=\lg x+\lg 3\) 的解是\fillinblank{}.
\begin{answer}
\(x=1\)，\(2\)
\end{answer}

\begin{solution}
对数有意义要求 \(x>0\)，此时 \(x^{2}+2>0\) 且 \(3>0\). 由对数的运算性质，原方程等价于
\(\lg\dfrac{x^{2}+2}{3x}=0\)，所以 \(\dfrac{x^{2}+2}{3x}=1\). 因为 \(x>0\)，分母不为零，整理得 \(x^{2}-3x+2=0\)，即 \((x-1)(x-2)=0\). 两个根 \(x=1\)，\(x=2\) 都满足定义域条件，故方程的解为 \(x=1\)，\(2\).
\end{solution}

\end{problem}

\begin{problem}
某地球仪上北纬 \(30^{\circ}\) 纬线的长度为 \(12\pi\mathrm{~cm}\)，该地球仪的半径是\fillinblank{}\(\mathrm{cm}\)，表面积是\fillinblank{}\(\mathrm{cm}^{2}\).
\begin{answer}
\(4\sqrt{3}\)，\(192\pi\)
\end{answer}

\begin{solution}
设地球仪半径为 \(R\). 北纬 \(30^{\circ}\) 纬线所在圆的半径是球心到该纬线平面的垂线在该平面内的投影长度，即 \(R\cos30^{\circ}=\dfrac{\sqrt3}{2}R\). 因此该纬线的长度为
\[
2\pi\cdot\frac{\sqrt3}{2}R=12\pi
\]
约去 \(\pi\) 并解得 \(R=4\sqrt3\). 地球仪的表面积为
\(4\pi R^{2}=4\pi\cdot(4\sqrt3)^{2}=192\pi\). 两个填空依次为 \(4\sqrt3\)，\(192\pi\).
\end{solution}

\end{problem}

\begin{problem}
圆 \(x^{2}+(y+1)^{2}=1\) 的圆心坐标是\fillinblank{}，如果直线 \(x+y+a=0\) 与该圆有公共点，那么实数 \(a\) 的取值范围是\fillinblank{}.
\begin{answer}
\((0,-1)\)，\([1-\sqrt{2},1+\sqrt{2}]\)
\end{answer}

\begin{solution}
将圆方程配方可知圆心为 \((0,-1)\)，半径为 \(1\). 圆心到直线 \(x+y+a=0\) 的距离为
\[
\frac{|0-1+a|}{\sqrt{1^{2}+1^{2}}}=\frac{|a-1|}{\sqrt2}
\]
直线与圆有公共点，当且仅当这个距离不大于半径 \(1\)，所以 \(\dfrac{|a-1|}{\sqrt2}\leqslant1\)，即 \(|a-1|\leqslant\sqrt2\). 因此 \(a\in[1-\sqrt2,1+\sqrt2]\).
\end{solution}

\end{problem}

\begin{problem}
在函数 \(f(x)=ax^{2}+bx+c\) 中，若 \(a\)，\(b\)，\(c\) 成等比数列且 \(f(0)=-4\)，则 \(f(x)\) 有\fillinblank{}值（填“大”或“小”），且该值为\fillinblank{}.
\begin{answer}
大，\(-3\)
\end{answer}

\begin{solution}
设等比数列 \(a\)，\(b\)，\(c\) 的公比为 \(q\). 因为 \(c=f(0)=-4\neq0\)，所以 \(a\neq0\)、\(q\neq0\)，且 \(c=aq^{2}\)，从而 \(a=-4/q^{2}<0\). 等比数列的相邻三项满足 \(b^{2}=ac\)，故
\[
f(x)=ax^{2}+bx+c
=a\left(x+\frac{b}{2a}\right)^{2}+c-\frac{b^{2}}{4a}
\]
由于 \(a<0\)，函数有最大值，且 \(c-\dfrac{b^{2}}{4a}=c-\dfrac{ac}{4a}=\dfrac34c=-3\). 所以填“大”，该值为 \(-3\).
\end{solution}

\end{problem}

\begin{problem}
定义“等和数列”：在一个数列中，如果每一项与它的后一项的和都为同一个常数，那么这个数列叫做等和数列，这个常数叫做该数列的公和. 已知数列 \(\{a_{n}\}\) 是等和数列，且 \(a_{1}=2\)，公和为 \(5\)，那么 \(a_{18}\) 的值为\fillinblank{}，且这个数列的前 \(21\) 项和 \(S_{21}\) 的计算公式为\fillinblank{}.
\begin{answer}
\(3\)，\(S_{21}=52\)
\end{answer}

\begin{solution}
由等和数列的定义，任意相邻两项满足 \(a_n+a_{n+1}=5\)，所以 \(a_{n+1}=5-a_n\). 由 \(a_1=2\) 得 \(a_2=3\)，再由递推关系得到 \(a_3=2\)，\(a_4=3\)，故可归纳得 \(a_{2j-1}=2\)，\(a_{2j}=3\)（\(j=1\)，\(2\)，\(\cdots\)）. 因为 \(18=2\times9\)，所以 \(a_{18}=3\). 前 \(21\) 项中有 \(11\) 个奇数项和 \(10\) 个偶数项，故
\(S_{21}=11\times2+10\times3=52\). 所以两个填空依次为 \(3\)，\(S_{21}=52\).
\end{solution}

\end{problem}

\section{解答题}
\begin{problem}
在 \(\triangle ABC\) 中，\(\sin A+\cos A=\frac{\sqrt{2}}{2}\)，\(AC=2\)，\(AB=3\)，求 \(\tan A\) 的值和 \(\triangle ABC\) 的面积.
\end{problem}

\begin{solution}
\begin{enumerate}
\item
因为 \(A\) 是三角形的内角，所以 \(0<A<\pi\). 由
\(\sin A+\cos A=\sqrt{2}\sin\left(A+\frac{\pi}{4}\right)\)，得
\(\sin\left(A+\frac{\pi}{4}\right)=\frac{1}{2}\). 令
\(t=A+\frac{\pi}{4}\)，则 \(t\in\left(\frac{\pi}{4},\frac{5\pi}{4}\right)\)，而
\(\sin t=\frac{1}{2}\) 的一般解为
\(t=\frac{\pi}{6}+2k\pi\) 或 \(t=\frac{5\pi}{6}+2k\pi\)（\(k\in\Z\)）. 因此在上述区间内只能取
\(t=\frac{5\pi}{6}\)，即 \(A=\frac{7\pi}{12}\).

于是 \(\tan A=\tan\left(\frac{\pi}{3}+\frac{\pi}{4}\right)=-(2+\sqrt{3})\)，且
\(\sin A=\sin\frac{7\pi}{12}=\frac{\sqrt{6}+\sqrt{2}}{4}\).

\item
因此
\[
S_{\triangle ABC}=\frac{1}{2}AB\cdot AC\sin A
=\frac{1}{2}\cdot3\cdot2\cdot\frac{\sqrt{6}+\sqrt{2}}{4}
=\frac{3(\sqrt{6}+\sqrt{2})}{4}
\]

故 \(\tan A=-(2+\sqrt{3})\)，\(\triangle ABC\) 的面积为
\(\dfrac{3(\sqrt{6}+\sqrt{2})}{4}\).
\end{enumerate}
\end{solution}

\begin{problem}
如图，在正三棱柱 \(ABC-A_{1}B_{1}C_{1}\) 中，\(AB=2\)，\(AA_{1}=2\)，由顶点 \(B\) 沿棱柱侧面经过棱 \(AA_{1}\) 到顶点 \(C_{1}\) 的最短路线与 \(AA_{1}\) 的交点为 \(M\)，求：

\begin{enumerate}
\item 该三棱柱的侧面展开图的对角线长；
\item 该最短路线的长与 \(\frac{A_{1}M}{AM}\) 的值；
\item 平面 \(C_{1}MB\) 与平面 \(ABC\) 所成二面角（锐角）的大小.
\end{enumerate}

\bitmapfigure[max width=0.42\linewidth]{img/2004/beijing_liberal_autumn/q16_fig1.png}
\end{problem}

\begin{solution}
\begin{enumerate}
\item 正三棱柱的三个侧面展开后组成一个长为
\(3AB=6\)，宽为 \(AA_{1}=2\) 的矩形，所以其对角线长为
\(\sqrt{6^{2}+2^{2}}=2\sqrt{10}\).

\item 把侧面 \(ABB_{1}A_{1}\) 绕棱 \(AA_{1}\) 展开到侧面
\(AA_{1}C_{1}C\) 所在平面，记点 \(B\) 的像为 \(D\). 在展开图中，取
\(A=(0,0)\)，\(A_{1}=(0,2)\)，则可取
\(D=(-2,0)\)，\(C_{1}=(2,2)\). 因而由两点间距离公式，最短路线的长度为
\(DC_{1}=\sqrt{(2-(-2))^{2}+(2-0)^{2}}=2\sqrt{5}\).

直线 \(DC_{1}\) 的方程为 \(y=\frac{1}{2}x+1\). 它与
\(AA_{1}\)（即 \(x=0\)）交于 \(M=(0,1)\)，所以 \(AM=A_{1}M=1\)，从而
\(\frac{A_{1}M}{AM}=1\).

\item 建立空间直角坐标系，取 \(A=(0,0,0)\)，\(B=(2,0,0)\)，\(C=(1,\sqrt{3},0)\)，
并取 \(A_{1}=(0,0,2)\)，\(C_{1}=(1,\sqrt{3},2)\)，\(M=(0,0,1)\).
平面 \(ABC\) 的一个法向量为 \(\bs{n}_{1}=(0,0,1)\). 平面
\(C_{1}MB\) 内的两个相交向量可取
\(\overrightarrow{MB}=(2,0,-1)\)，\(\overrightarrow{MC_{1}}=(1,\sqrt{3},1)\)，所以它的一个法向量为
\[
\bs{n}_{2}=\overrightarrow{MB}\times\overrightarrow{MC_{1}}
=(\sqrt{3},-3,2\sqrt{3})
\]
于是 \(\left\lvert\bs{n}_{2}\right\rvert=2\sqrt{6}\)，
\(\left\lvert\bs{n}_{1}\cdot\bs{n}_{2}\right\rvert=2\sqrt{3}\).
设两平面的锐二面角为 \(\theta\)，则
\[
\cos\theta=\frac{\left\lvert\bs{n}_{1}\cdot\bs{n}_{2}\right\rvert}
{\left\lvert\bs{n}_{1}\right\rvert\left\lvert\bs{n}_{2}\right\rvert}
=\frac{2\sqrt{3}}{2\sqrt{6}}=\frac{\sqrt{2}}{2}
\]
因此 \(\theta=45^{\circ}\).
\end{enumerate}
\end{solution}

\begin{problem}
如图，抛物线关于 \(x\) 轴对称，它的顶点在坐标原点，点 \(P(1,2)\)，\(A(x_{1},y_{1})\)，\(B(x_{2},y_{2})\) 均在抛物线上.

\begin{enumerate}
\item 写出该抛物线的方程及其准线方程；
\item 当 \(PA\) 与 \(PB\) 的斜率存在且倾斜角互补时，求 \(y_{1}+y_{2}\) 的值及直线 \(AB\) 的斜率.
\end{enumerate}

\bitmapfigure[max width=0.58\linewidth]{img/2004/beijing_liberal_autumn/q17_fig1.png}
\end{problem}

\begin{solution}
\begin{enumerate}
\item 抛物线的顶点在原点且关于 \(x\) 轴对称，结合图形可设其方程为
\(y^{2}=2px\)，其中 \(p>0\). 因为点 \(P(1,2)\) 在抛物线上，所以
\(2^{2}=2p\cdot1\)，解得 \(p=2\). 因此抛物线方程为 \(y^{2}=4x\)，焦点为 \((1,0)\)，准线方程为
\(x=-1\).

\item 由 \(A\)，\(B\) 在抛物线上，得
\(x_{1}=\dfrac{y_{1}^{2}}{4}\)，\(x_{2}=\dfrac{y_{2}^{2}}{4}\). 又因直线
\(PA\)，\(PB\) 的斜率存在，所以 \(y_{1}\neq2\)，\(-2\)，\(y_{2}\neq2\)，\(-2\). 因此
\[
k_{PA}=\frac{y_{1}-2}{x_{1}-1}
=\frac{y_{1}-2}{\dfrac{y_{1}^{2}}{4}-1}
=\frac{4}{y_{1}+2}
\]
且 \(k_{PB}=\frac{4}{y_{2}+2}\).
设直线 \(PA\)，\(PB\) 的倾斜角分别为 \(\alpha\)，\(\beta\). 因为两角互补，
\(\alpha+\beta=\pi\)，所以
\(k_{PA}=\tan\alpha=-\tan\beta=-k_{PB}\). 从而
\(\frac{4}{y_{1}+2}=-\frac{4}{y_{2}+2}\)，即 \(y_{1}+y_{2}=-4\).

由 \(y_{1}+y_{2}=-4\) 可知 \(y_{1}\neq y_{2}\)，否则将有 \(y_{1}=y_{2}=-2\)，这与 \(PA\) 的斜率存在矛盾. 再由
\(x_{2}-x_{1}=\frac{y_{2}^{2}-y_{1}^{2}}{4}=\frac{(y_{2}-y_{1})(y_{1}+y_{2})}{4}\)，可得
\(k_{AB}=\frac{y_{2}-y_{1}}{x_{2}-x_{1}}=\frac{4}{y_{1}+y_{2}}=-1\).
故 \(y_{1}+y_{2}=-4\)，直线 \(AB\) 的斜率为 \(-1\).
\end{enumerate}
\end{solution}

\begin{problem}
函数 \(f(x)\) 是定义在 \([0,1]\) 上的增函数，满足 \(f(x)=2f\left(\frac{x}{2}\right)\) 且 \(f(1)=1\)，在每个区间 \(\left(\frac{1}{2^{i}},\frac{1}{2^{i-1}}\right]\)（\(i=1\)，\(2\)，\(\cdots\)）上，\(y=f(x)\) 的图象都是斜率为同一常数 \(k\) 的直线的一部分.

\begin{enumerate}
\item 求 \(f(0)\) 及 \(f\left(\frac{1}{2}\right)\)，\(f\left(\frac{1}{4}\right)\) 的值，并归纳出 \(f\left(\frac{1}{2^{i}}\right)\)（\(i=1\)，\(2\)，\(\cdots\)）的表达式；
\item 设直线 \(x=\frac{1}{2^{i}}\)，\(x=\frac{1}{2^{i-1}}\)，\(x\) 轴及 \(y=f(x)\) 的图象围成的矩形的面积为 \(a_{i}\)（\(i=1\)，\(2\)，\(\cdots\)），求 \(a_{1}\)，\(a_{2}\) 及 \(\displaystyle\lim_{n\to\infty}(a_{1}+a_{2}+\cdots+a_{n})\) 的值.
\end{enumerate}
\end{problem}

\begin{solution}
\begin{enumerate}
\item 由 \(f(0)=2f(0)\)，得 \(f(0)=0\). 又由
\(f(1)=2f\left(\frac{1}{2}\right)=1\)，得
\(f\left(\frac{1}{2}\right)=\frac{1}{2}\). 再由
\(f\left(\frac{1}{2}\right)=2f\left(\frac{1}{4}\right)\)，得
\(f\left(\frac{1}{4}\right)=\frac{1}{4}\). 依此类推，得到
\(f\left(\frac{1}{2^{i}}\right)=\frac{1}{2^{i}}\)（\(i=1\)，\(2\)，\(\cdots\)）.

\item
由 \(f\) 在 \([0,1]\) 上递增且 \(f(0)=0\)，有 \(f(x)\geqslant0\). 每个区间上的图象是一条斜率为同一常数 \(k\) 的线段. 题设还要求它与两条竖直边、\(x\) 轴围成矩形；矩形的上边必须平行于 \(x\) 轴，所以该线段必须水平，故公共斜率 \(k=0\). 对于
\(\frac{1}{2^{i}}<x\leqslant\frac{1}{2^{i-1}}\)，右端点属于该区间，且图象水平，因此
\(f(x)=f\left(\frac{1}{2^{i-1}}\right)=\frac{1}{2^{i-1}}\).

矩形的宽为
\(\frac{1}{2^{i-1}}-\frac{1}{2^{i}}=\frac{1}{2^{i}}\)，故
\(a_{i}=\frac{1}{2^{i-1}}\cdot\frac{1}{2^{i}}=\frac{1}{2^{2i-1}}\).
于是 \(a_{1}=\frac{1}{2}\)，\(a_{2}=\frac{1}{8}\)，
且
\[
\lim_{n\to\infty}(a_{1}+a_{2}+\cdots+a_{n})
=\frac{\frac{1}{2}}{1-\frac{1}{4}}=\frac{2}{3}
\]
\end{enumerate}
\end{solution}

\begin{problem}
某段城铁线路上依次有 \(A\)，\(B\)，\(C\) 三站，\(AB=15\mathrm{~km}\)，\(BC=3\mathrm{~km}\)，在列车运行时刻表上，规定列车 8 时整从 \(A\) 站发车，8 时 07 分到达 \(B\) 站并停车 1 分钟，8 时 12 分到达 \(C\) 站，在实际运行中，假设列车从 \(A\) 站正点发车，在 \(B\) 站停留 1 分钟，并在行驶时以同一速度 \(v\mathrm{~km}/\mathrm{h}\) 匀速行驶，列车从 \(A\) 站到达某站的时间与时刻表上相应时间之差的绝对值称为列车在该站的运行误差.

\begin{enumerate}
\item 分别写出列车在 \(B\)，\(C\) 两站的运行误差；
\item 若要求列车在 \(B\)，\(C\) 两站的运行误差之和不超过 2 分钟，求 \(v\) 的取值范围.
\end{enumerate}
\end{problem}

\begin{solution}
\begin{enumerate}
\item 以分钟为单位计算运行误差. 速度满足 \(v>0\). 列车从 \(A\) 到 \(B\) 的实际运行时间为
\(\dfrac{15}{v}\) 小时，即 \(\dfrac{900}{v}\) 分钟. 时刻表规定从 8 时整到 8 时 07 分经过 \(7\) 分钟，所以
\(e_{B}=\left\lvert\frac{900}{v}-7\right\rvert\).
列车从 \(A\) 到 \(C\) 的实际行驶时间为 \(\dfrac{18}{v}\) 小时，另加在 \(B\) 站停留的 \(1\) 分钟，因此实际到达 \(C\) 站的时刻距 \(8\) 时
\(\dfrac{1080}{v}+1\) 分钟，故
\(e_{C}=\left\lvert\frac{1080}{v}+1-12\right\rvert=\left\lvert\frac{1080}{v}-11\right\rvert\).

\item 令 \(t=\dfrac{180}{v}>0\)，则条件 \(e_{B}+e_{C}\leqslant2\) 等价于
\(\lvert5t-7\rvert+\lvert6t-11\rvert\leqslant2\).
分三种情况讨论，两个绝对值内的式子分别在 \(t=\frac{7}{5}\) 和 \(t=\frac{11}{6}\) 处变号：

当 \(0<t\leqslant\frac{7}{5}\) 时，左端为
\(18-11t\geqslant\frac{13}{5}>2\)；

当 \(\frac{7}{5}\leqslant t\leqslant\frac{11}{6}\) 时，左端为
\(4-t\geqslant\frac{13}{6}>2\)；

当 \(t\geqslant\frac{11}{6}\) 时，左端为
\(11t-18\geqslant\frac{13}{6}>2\).

因此对任意 \(t>0\)，都有
\(\lvert5t-7\rvert+\lvert6t-11\rvert\geqslant\frac{13}{6}>2\)，不存在满足条件的正数 \(v\). 故
\(v\) 的取值范围为 \(\varnothing\).
\end{enumerate}
\end{solution}

\begin{problem}
给定有限个正数满足条件 \(T\)：每个数都不大于 \(50\) 且总和 \(L=1275\). 现将这些数按下列要求进行分组，每组数之和不大于 \(150\) 且分组的步骤是：首先，从这些数中选择这样一些数构成第一组，使得 \(150\) 与这组数之和的差 \(r_{1}\) 与所有可能的其他选择相比是最小的，\(r_{1}\) 称为第一组余差；然后，在去掉已选入第一组的数后，对余下的数按第一组的选择方式构成第二组，这时的余差为 \(r_{2}\)；如此继续构成第三组（余差为 \(r_{3}\)），第四组（余差为 \(r_{4}\)），\(\cdots\)，直至第 \(N\) 组（余差为 \(r_{N}\)）把这些数全部分完为止.

\begin{enumerate}
\item 判断 \(r_{1}\)，\(r_{2}\)，\(\cdots\)，\(r_{N}\) 的大小关系，并指出除第 \(N\) 组外的每组至少含有几个数；
\item 当构成第 \(n\)（\(n<N\)）组后，指出余下的每个数与 \(r_{n}\) 的大小关系，并证明 \(r_{n-1}>\frac{150n-L}{n-1}\)；
\item 对任何满足条件 \(T\) 的有限个正数，证明：\(N\leqslant 11\).
\end{enumerate}
\end{problem}

\begin{solution}
\begin{enumerate}
\item 设第 \(i\) 组形成前剩余的数构成集合 \(E_i\)，记该组所取数的和为
\(150-r_i\). 由于每组的和不大于 \(150\)，有 \(r_i\geqslant0\). 形成第 \(i+1\) 组时，可选的数只来自去掉第 \(i\) 组后的集合，因此所有可选组合都属于前一步的可选组合，能取得的最大组和不增加，故
\[
r_{1}\leqslant r_{2}\leqslant\cdots\leqslant r_{N}
\]

若第 \(i\) 组不是最后一组而只含有至多两个数，则该组的和不大于
\(2\times50=100\)，从而 \(r_i\geqslant50\). 但第 \(i\) 组之后仍有一个正数，设为 \(x\)，则
\(0<x\leqslant50\leqslant r_i\)，把 \(x\) 加入第 \(i\) 组后组和仍不大于 \(150\)，且余差变小，这与 \(r_i\) 的最小性矛盾. 所以除第 \(N\) 组外，每组至少含有 \(3\) 个数.
\item 第 \(n\) 组形成后仍有数未分组. 若余下某个数 \(x\leqslant r_n\)，则第 \(n\) 组加入该数后的和为 \(150-r_n+x\leqslant150\)，且因 \(x>0\) 其余差为 \(r_n-x<r_n\)，这与 \(r_n\) 的最小性矛盾. 因而余下的每个数都满足
\[
x>r_n
\]

余下至少有一个数，故余下数的总和大于 \(r_n\). 而前 \(n\) 组的总和为
\(\sum_{i=1}^{n}(150-r_i)\)，所以
\[
L-\sum_{i=1}^{n}(150-r_i)>r_n
\]
整理得
\[
r_1+r_2+\cdots+r_{n-1}>150n-L
\]
对 \(2\leqslant n<N\)，由（1）知
\(r_i\leqslant r_{n-1}\)（\(1\leqslant i\leqslant n-1\)），于是
\[
(n-1)r_{n-1}\geqslant r_1+r_2+\cdots+r_{n-1}>150n-L
\]
除以正数 \(n-1\)，因此
\[
r_{n-1}>\frac{150n-L}{n-1}
\]

\item 反设 \(N>11\). 由（2）取 \(n=11\)，有
\[
r_{10}>\frac{150\times11-1275}{10}=37.5
\]
由于 \(r_{11}\geqslant r_{10}\)，所以 \(r_{11}>37.5\). 又因第 \(11\) 组不是最后一组，它至少含有三个数；第 \(11\) 组中的数是在形成第 \(10\) 组后剩下的数，由（2）中已经证明的结论（取 \(n=10\)）可知每个都大于 \(r_{10}>37.5\). 设第 \(11\) 组的和为 \(s_{11}\)，则

\[
s_{11}>3\times37.5=112.5
\]

从而
\[
r_{11}=150-s_{11}<37.5
\]
这与 \(r_{11}>37.5\) 矛盾. 故假设不成立，必有
\[
N\leqslant11
\]
\end{enumerate}
\end{solution}
